% O014 / D80 — Methods of Algebra, Volume 2 — id-ID
% Sumber: appendix1.tex, baris 488--531 (inklusif)
% ID stabil unit: o014.aljabr2.appendix-a.exercises
\begin{Exercises}
	\item Dalam pembahasan pada \S\ref{sec:Yoneda}, tetapkan suatu gelanggang komutatif $\Bbbk$, syaratkan bahwa $\mathcal{C}$ adalah kategori-$\Bbbk\dcate{Mod}$ (Definisi \ref{def:cat-k-linear}), gantikan $\cate{Set}$ dengan $\Bbbk\dcate{Mod}$ dalam definisi $\mathcal{C}^\wedge$ dan $\mathcal{C}^\vee$, serta syaratkan semua funktor bersifat $\Bbbk$-linear. Rumuskan versi $\Bbbk$-linear yang bersesuaian dari Teorema \ref{prop:Yoneda} dan \ref{prop:Yoneda-density}.
	
	\item Lengkapi arah ``hanya jika'' pada Contoh \ref{eg:Mod-cpt}.
	\begin{hint}
		Setiap modul-$R$ dapat ditulis sebagai $\varinjlim$ terfilter dari modul-$R$ berpresentasi hingga (yang tidak harus berupa submodulnya); demi kesederhanaan tulislah $X = \varinjlim \alpha$. Perhatikan prapeta dari $\identity_X$ melalui \eqref{eqn:I-small-gen}; diperoleh $f_i \in \Hom(X, \alpha(i))$ sedemikian sehingga $\iota_i f_i = \identity_X$. Ini memberikan dekomposisi $\alpha(i) \simeq X \oplus N$. Karena $N$ isomorfik dengan suatu hasil bagi dari $\alpha(i)$, ia merupakan modul-$R$ yang dibangun secara hingga. Masalahnya kini tereduksi pada fakta bahwa hasil bagi modul-$R$ berpresentasi hingga oleh submodul yang dibangun secara hingga tetap berpresentasi hingga.
	\end{hint}
	
	\item Jelaskan hubungan antara ordinal kofinal yang diperkenalkan pada \S\ref{sec:cpt-objects} dan subkategori kofinal pada Definisi \ref{def:cofinal}.
	
	\item (J.\ Adámek, J.\ Rosický) Buktikan fakta yang disebutkan dalam Catatan \ref{rem:filtered-poset} dengan cara berikut: untuk setiap kategori kecil terfilter $I$, terdapat suatu himpunan kecil terurut parsial terfilter $(\mathcal{P}, \leq)$ beserta funktor kofinal $\mathcal{P} \to I$.
	\begin{enumerate}[(i)]
		\item Andaikan kategori kecil terfilter $I$ mempunyai sifat berikut: setiap subkategori hingga $\mathcal{A}$ dapat disertakan ke dalam suatu subkategori hingga $\mathcal{A}'$ yang mempunyai objek terminal tunggal (perhatikan bahwa subkategori tidak harus penuh, sedangkan kategori hingga berarti bahwa himpunan objek dan himpunan morfismenya sama-sama hingga). Misalkan $\mathcal{I}$ adalah himpunan semua subkategori yang mempunyai objek terminal tunggal, dengan urutan parsial oleh inklusi $\subset$. Tunjukkan bahwa dapat didefinisikan sebuah funktor
		\[ H: \mathcal{I} \to I, \]
		yang memetakan objek $\mathcal{A}$ ke objek terminal tunggalnya $H\mathcal{A}$ dan memetakan morfisme $\mathcal{A}_1 \subset \mathcal{A}_2$ ke morfisme tunggal $H\mathcal{A}_1 \to H\mathcal{A}_2$ di dalam $\mathcal{A}_2$.
		\item Dengan notasi di atas, tunjukkan bahwa $(\mathcal{I}, \subset)$ adalah himpunan kecil terurut parsial terfilter dan bahwa $H$ adalah funktor kofinal.
		\item Sekarang misalkan $I$ adalah kategori kecil terfilter sebarang. Tunjukkan bahwa kategori produk $I \times \omega$ mempunyai sifat yang diasumsikan dalam (i), sedangkan funktor proyeksi $I \times \omega \to I$ bersifat kofinal. Di sini $\omega$ adalah himpunan terurut total bilangan bulat taknegatif.
		
		\begin{hint}
			Jelas bahwa $I \times \omega$ terfilter. Untuk suatu subkategori hingga $\mathcal{A}$ dari $I \times \omega$, tunjukkan bahwa terdapat objek $(i, n)$ dari $I \times \omega$ beserta sebuah kerucut dengan funktor inklusi $\mathcal{A} \to I \times \omega$ sebagai alas dan $(i, n)$ sebagai puncaknya (Konvensi \ref{con:limit-cone}). Tambahkan objek $(i, n+1)$ kepada $\mathcal{A}$, beserta morfisme identitasnya dan semua komposisi dari morfisme $(j, k) \to (i, n)$ dalam kerucut dengan morfisme kanonik $(i, n) \to (i, n+1)$. Dengan demikian diperoleh subkategori hingga $\mathcal{A}'$ yang mempunyai objek terminal tunggal $(i, n+1)$.
		\end{hint}
	
		\item Gunakan (iii) untuk membuktikan keberadaan funktor kofinal $\mathcal{P} \to I$ yang diminta.
	\end{enumerate}
	
	\item Misalkan kategori $\mathcal{C}$ terakses.
	\begin{enumerate}[(i)]
		\item Pilih himpunan kecil $S \subset \Obj(\mathcal{C})$ seperti dalam Definisi \ref{def:presentable-cat}, dan nyatakan dengan $\mathcal{S}$ subkategori penuh yang dibangunnya. Buktikan bahwa pembenaman Yoneda membatasi menjadi funktor setia-penuh $\mathcal{C} \to \cate{Set}^{\mathcal{S}^{\opp}}$.
		\item Buktikan bahwa $\mathcal{C}$ berdaya-baik (Definisi \ref{def:well-powered}).
		\begin{hint}
			Untuk setiap objek $F$ dari $\cate{Set}^{\mathcal{S}^{\opp}}$ dan $X \in S$, himpunan $FX$ dan himpunan kuasanya sama-sama kecil. Gunakan fakta ini untuk menunjukkan bahwa $\mathrm{Sub}_F$ adalah himpunan kecil.
		\end{hint}
	\end{enumerate}

	% Rujukan: Deligne, Categories tannakiennes, 2.18
	\item Misalkan $\Bbbk$ sebuah medan. Misalkan $\mathcal{A}$ adalah kategori abelian $\Bbbk$-linear yang hingga secara lokal, sedangkan $\mathcal{A}'$ adalah subkategori abeliannya, dan tuliskan funktor inklusinya sebagai $i: \mathcal{A}' \to \mathcal{A}$. Andaikan subhimpunan $\Obj(\mathcal{A}') \subset \Obj(\mathcal{A})$ tertutup terhadap isomorfisme.
	\begin{enumerate}[(i)]
		\item Tunjukkan bahwa untuk setiap $X \in \Obj(\mathcal{A})$, irisan poset objek hasil bagi $\mathrm{Quot}_X$ (atau poset subobjek $\mathrm{Sub}_X$) dengan $\mathcal{A}'$ mempunyai unsur terbesar tunggal, yang dinyatakan dengan $i^*(X)$ (atau $i^!(X)$), dan bahwa $i^*$ (atau $i^!$) memberikan funktor adjoin kiri (atau kanan) dari $i$.
		\item Misalkan $\mathcal{A}$ mempunyai pembangkit projektif $s$. Buktikan bahwa $s' := i^*(s)$ adalah pembangkit projektif bagi $\mathcal{A}'$.
		\item Untuk $s,s'$ seperti di atas, misalkan $A := \End(s)$ dan $A' := \End(s')$. Buktikan bahwa
		\[ A = \Hom(s, s) \to \Hom(s, s') \leftiso \Hom(s', s') = A' \]
		menjadikan $A'$ gelanggang hasil bagi dari $A$; tuliskan $A' = A/\mathfrak{a}$. Buktikan bahwa $\mathcal{A}$ dan $s$ memenuhi syarat-syarat Teorema \ref{prop:identify-ModR-fg}, sehingga $\mathcal{A}$ ekuivalen dengan $\catesubd{Mod}{\mathrm{fg}}A$, sedangkan $\mathcal{A}'$ ekuivalen dengan subkategori penuh yang terdiri atas modul-modul yang dianihilasi oleh ideal dua-sisi $\mathfrak{a}$.
	\end{enumerate}
\end{Exercises}
